% Hướng dẫn cho Phản biện. Dịch: cd docs/huong-dan/pb && pdflatex slides.tex && pdflatex slides.tex
% (trước đó chạy scripts/test-all.sh để có ảnh trong out/ui/)
\documentclass[aspectratio=169,12pt]{beamer}
\usepackage{../pi-hd}
\newcommand\vaitro{Hướng dẫn cho Phản biện}
\begin{document}

\bia{Phản biện}{Đọc bài được giao, tự giải, điền phiếu phản biện, trao đổi ẩn danh.}

\chuanbi
\dungthu

\chu{Thầy cô làm gì, trong bao lâu}{%
\begin{buoc}
\item Nhận \textbf{thư mời} (người gửi: „Pi — Đề ra kỳ này"): số bài, hạn phản biện, đường dẫn đăng nhập.
\item Mở từng bài, \textbf{tự giải trước} — lời giải chỉ hiện khi Phụ trách chuyên mục mở cho cả kỳ.
\item Điền \textbf{phiếu phản biện} cho mỗi bài: mức đề nghị, đề nghị chọn hay không, nhận xét.
\item \textbf{Đánh dấu đã xong}. Trao đổi thêm ở mục Thảo luận của bài nếu cần.
\end{buoc}
\medskip
Chưa xong thì có thư nhắc 3 ngày và 1 ngày trước hạn. Hết kỳ, các bài rời khỏi danh sách của thầy cô.\\[2mm]
\textbf{Chấm ẩn danh}: thầy cô không thấy tên tác giả; phản biện khác thấy thầy cô là ``Phản biện 1, 2\ldots'', không thấy email.}

\manhinh[Mỗi lần vào hệ thống đều đi qua trang này.]{dang-nhap}{Vào hệ thống}
\manhinh{pb-danh-sach}{Danh sách: chỉ các bài được giao}
\manhinh{pb-phieu}{Trang bài và phiếu phản biện}
\manhinh{pb2-an-danh}{Thảo luận ẩn danh}
\manhinh{pb-loi-giai}{Khi lời giải được mở}
\anh{dien-thoai-bai}{Trên điện thoại}{Trang dùng được trên điện thoại: đọc đề, điền phiếu, thảo luận.\\[2mm]
Công thức dài cuộn ngang trong khung riêng.\\[2mm]
Công thức trong nhận xét viết như \LaTeX: \texttt{\$x>0\$} hiện thành $x>0$.}

\chu{Khi có sự cố}{%
\begin{itemize}\setlength\itemsep{2mm}
\item \textbf{„Phiên vào hệ thống đã hết hạn"}: bấm \emph{Đăng nhập lại} (liên kết có hiệu lực 10 phút; phiên làm việc 6 giờ).
\item \textbf{Dòng „Xin chào …" không phải tài khoản của thầy cô}: trình duyệt đang đăng nhập nhiều tài khoản Google —
  mở cửa sổ riêng tư, chỉ đăng nhập đúng một tài khoản.
\item \textbf{Không thấy bài nào}: kỳ phản biện đã đóng, hoặc chưa được giao bài — hỏi Phụ trách chuyên mục.
\item \textbf{Bấm „Lưu phiếu" bị nhắc}: chưa chọn đề nghị (chọn / sửa rồi chọn / không chọn).
\item \textbf{Muốn sửa phiếu đã đánh dấu xong}: bấm \emph{Bỏ đánh dấu xong}, sửa, lưu, đánh dấu lại.
\end{itemize}}

\gopy

\chu{Tóm tắt}{%
\begin{buoc}
\item Mở đường dẫn trong thư mời, kiểm tra dòng „Xin chào", bấm \emph{Vào hệ thống}.
\item Mở từng bài, tự giải.
\item Điền phiếu: mức, đề nghị, nhận xét — \emph{Lưu phiếu}.
\item \emph{Đánh dấu đã xong} trước hạn.
\item Trao đổi ở Thảo luận; khi lời giải được mở, đối chiếu và nhận xét thêm.
\end{buoc}}

\end{document}
